% ======================================================================
\section{Classical nonlinear manifold learning}
\label{sec:manifold}
% ======================================================================

Before neural autoencoders became practical, a family of nonlinear DR methods
was developed in machine learning to recover low-dimensional coordinates for
data lying near curved manifolds. These methods work directly with the data
set, often through a similarity graph or kernel matrix, and learn an
\emph{embedding} of the training points. Most do not provide a decoder, which
has limited their use as reduced representations for dynamics but made them
valuable for analysis and for estimating intrinsic dimension.

\subsection{Kernel principal component analysis}

Kernel PCA \parencite{scholkopf1998kpca} performs PCA in a feature space defined
implicitly by a positive-definite kernel $k(\bm{x},\bm{x}')$, for example a
Gaussian or polynomial kernel. The embedding coordinates follow from the
eigenvectors $\bm{\alpha}_j$ of the centered $T\times T$ kernel matrix,
$z_j(\bm{x})=\sum_{i}\alpha_{j,i}\,k(\bm{x}_i,\bm{x})$, so kernel PCA is the
nonlinear analogue of the method of snapshots and has the same cost scaling.
Reconstruction requires solving a \emph{pre-image} problem, finding a state
whose feature-space image is close to a given point, which is ill-posed and
usually solved approximately by fixed-point iteration or regression
\parencite{mika1999preimage}. \textcite{salvador2021kpod} used kernel POD to
build reduced spaces for parametrized PDEs and trained neural networks to
predict the reduced coefficients. They obtained more accurate snapshot
approximations with fewer coordinates than POD at similar online cost.

\subsection{Graph-based spectral embeddings}

Isomap \parencite{tenenbaum2000isomap} approximates geodesic distances on the
manifold by shortest paths in a nearest-neighbour graph and embeds them with
multidimensional scaling. Locally linear embedding (LLE)
\parencite{roweis2000lle} computes weights that reconstruct each point from its
neighbours and finds low-dimensional coordinates that preserve those weights.
Laplacian eigenmaps \parencite{belkin2003laplacian} and diffusion maps
\parencite{coifman2006diffusion} use eigenvectors of a graph Laplacian or of a
Markov diffusion operator on the data. Diffusion maps add a time parameter
that exposes the geometry at multiple scales. All of these methods preserve
local neighbourhoods and can unfold manifolds that linear projections fold
onto themselves.

In fluid dynamics, \textcite{farzamnik2023snapshots} combined Isomap as an
encoder with a $k$-nearest-neighbour decoder. For the fluidic pinball they
described both the pitchfork bifurcation and the chaotic regime with only
three coordinates, linked to vortex-shedding phases and force coefficients.
For the wake of tandem cylinders, Isomap separated two shedding regimes that
appear as a single circle in the plane of the first two POD amplitudes.
\textcite{csala2022comparing} compared four manifold-learning and two
autoencoder methods with PCA on canonical laminar and turbulent flows and on
biomedical flows. They discussed the influence of arranging data temporally
versus spatially and concluded that nonlinear DR is promising for reduced-order
modeling. \textcite{mendez2023linear} placed kernel PCA, Isomap, and LLE in a
common framework with autoencoders and benchmarked them for filtering,
identification of oscillatory patterns, and compression.
\textcite{loiseau2018sparse} lifted sensor signals to a feature space by time
delays, identified sparse dynamics on it, and reconstructed the velocity field
with a local linear map from features to flow. This is another way to supply
the decoder that spectral embeddings lack.

The limitations of these methods for reduced-order modeling are well
recognized \parencite{mojgani2021registration}. There is no explicit map from
latent coordinates back to the physical space, and out-of-sample points need
Nystr\"om-type extensions. The embeddings depend on neighbourhood parameters,
and the coordinates are often hard to interpret. Cost and memory scale with
$T^{2}$. Visualization methods such as t-SNE \parencite{vandermaaten2008tsne}
and UMAP \parencite{mcinnes2018umap} are useful for exploring the structure of
flow data sets, but they deliberately distort global geometry and are not
suitable as coordinates for dynamics models. Manifold-learning estimates of
intrinsic dimension, such as the maximum-likelihood estimator of
\textcite{levina2004intrinsic}, remain useful for choosing latent sizes of the
autoencoders discussed next.

\subsection{Polynomial and quadratic manifolds}

A simple way to curve a linear subspace while keeping an explicit decoder is
to add polynomial corrections,
\begin{equation}
  \q\approx\bar{\q}+\bm{V}\z+\bm{H}\,(\z\otimes\z),
  \label{eq:quadratic}
\end{equation}
where $\bm{V}$ is usually a POD basis and $\bm{H}$ is fitted to the residual.
Quadratic manifolds were introduced for nonlinear structural dynamics
\parencite{jain2017quadratic}. They were later developed to mitigate the
Kolmogorov barrier in projection-based model reduction
\parencite{barnett2022quadratic} and combined with operator inference for
non-intrusive models \parencite{geelen2023quadratic}. The decoder
\eqref{eq:quadratic} is cheap and differentiable, the encoder can remain a
linear projection, and the construction is a small step from standard POD
workflows. The expressiveness is limited compared with deep networks, but
quadratic manifolds provide a useful intermediate baseline between linear and
fully learned trial manifolds.
